\section{Três dimensões}
\label{sec:tres-dimensoes}

A base tridimensional foi construída depois da bidimensional, em arquivos
separados, pela razão de projeto da seção~\ref{sec:estrategias}. O escopo foi
declarado antes de começar e cumprido sem alargamento: \textbf{só o que o solver
de Navier--Stokes vai usar} --- malha, ajuste por mínimos quadrados móveis,
estêncil, campos e o Poisson que os verifica. Compactos, laplaciano pontual e
lote em GPU ficaram de fora, declaradamente.

Esta seção registra o que \emph{mudou} de duas para três dimensões, e a parte
que importa não é trocar um 2 por um 3.

\subsection{O teto teórico que nunca é atingido}

Um hexaedro tem seis faces, e na transição 2:1 cada face pode ter
\textbf{quatro} vizinhos --- contra duas faces e dois vizinhos em duas dimensões.
O teto teórico é $6\times 4 = 24$ vizinhos por célula, e dimensionar a carga da
troca por ele é a escolha óbvia.

A medida (\texttt{bench/franja3d}, feita \emph{antes} de o código tridimensional
existir) a desautoriza: o máximo real é \textbf{9} em malha adaptada e 6 em
uniforme. Uma célula tem, no máximo, uma face com vizinhos refinados em cada
direção --- não as seis. Dimensionar pelo teto custaria \num{2.7} vezes mais
memória, e a classe de malha passou a \emph{medir} os vizinhos na construção em
vez de presumi-los.

\subsection{O suporte em cruz não sustenta ordem 3}

A base polinomial explode com a dimensão: 4, 10, 20 e 35 monômios para as ordens
1 a 4, contra 3, 6, 10 e 15 em duas dimensões. O suporte mínimo é o número de
monômios, de modo que a ordem 3 pede 20 pontos onde antes pedia 10.

Crescer o suporte por \textbf{anéis de vizinhos de face} --- a estratégia que
funciona em duas dimensões --- falha aqui, e falha quase sempre: numa malha
uniforme de nível 3, \textbf{504 dos 512} ajustes de ordem 3 rebaixaram a ordem,
\SI{98.4}{\percent}. Os pontos em cruz são colineares com o alvo em cada direção,
e colinearidade não sustenta monômios cruzados.

A substituição é o suporte por \textbf{caixa geométrica}, e o raio dela sai da
álgebra, não de escolha: uma caixa de raio $r$ tem $2r+1$ valores distintos por
direção, e reproduzir grau $g$ exige $g+1$ deles, logo $r \geq \lceil g/2\rceil$.
A primeira versão usava raio 1 --- 27 pontos para 20 monômios, número de sobra ---
e ainda rebaixava em 216 dos 512 ajustes. Os 216 eram exatamente $6^3$: as células
interiores, onde uma caixa $3\times3\times3$ tem só três valores distintos por
direção e $x^3 = h^2 x$ sobre eles. \textbf{Número de pontos suficiente não é
posicionamento suficiente}, e o contador de rebaixamento foi o que disse isso.

\subsection{Publicar a caixa, e não propagar por faces}

A célula mais distante da caixa de raio 2 está a seis faces de distância.
Propagando por faces, cobri-la exigia cinco camadas e \SI{200.7}{\kilo\byte} por
folha. Publicando a \emph{caixa}, o dono de um fantasma meu tem como locais
justamente as células que me faltam:

\begin{center}
\begin{tabular}{lrrr}
\toprule
camadas & 2 & 3 & 5 \\
\midrule
truncados, propagação topológica & 150 & 48 & 0 \\
truncados, publicação da caixa   &  69 & 0  & 0 \\
KB por folha                     & 4,3 & 19,6 & 114,0 \\
\bottomrule
\end{tabular}
\end{center}

Zero truncamento passou de cinco camadas e \SI{200.7}{\kilo\byte} para três
camadas e \SI{19.6}{\kilo\byte}: \textbf{dez vezes menos memória}.

\reserva{o primeiro resultado desta medida foi ``zero truncamento com qualquer
número de camadas'', que é bom demais. O suporte por caixa descartava em silêncio
os pontos que não achava, e as duas respostas negativas de uma busca geométrica
são diferentes: um ponto \emph{fora do domínio} é legítimo --- o contorno trata
---, e um ponto \emph{dentro} que não se acha significa que falta franja.
Confundi-las esconde a segunda. Os números acima são os de depois da distinção.}

\subsection{O Poisson em três dimensões}

O oráculo primário continua sendo o \textbf{grau de exatidão}, não a ordem de
convergência, e cada afirmação de exatidão vem com o contraste onde ela deve
falhar.

\begin{center}
\begin{tabular}{llrr}
\toprule
solução & esquema & malha & erro $\infty$ \\
\midrule
\texttt{grau1} & fluxo 2ª ordem    & uniforme, $8^3$   & \num{1.5447e-11} \\
\texttt{grau2} & fluxo 2ª ordem    & uniforme, $8^3$   & \num{2.9505e-3} \\
\texttt{grau2} & MLS+Gauss ordem 2 & uniforme, $8^3$   & \num{5.0082e-12} \\
\texttt{grau2} & fluxo 2ª ordem    & \textbf{adaptada} & \num{1.0451e-1} \\
\texttt{grau2} & MLS+Gauss ordem 2 & \textbf{adaptada} & \num{1.9738e-11} \\
\texttt{grau3} & MLS+Gauss ordem 3 & uniforme, $8^3$   & \num{1.2479e-13} \\
\texttt{grau4} & MLS+Gauss ordem 3 & uniforme, $8^3$   & \num{2.4062e-3} \\
\bottomrule
\end{tabular}
\end{center}

As linhas andam em pares e é o par que afirma algo. O fluxo de segunda ordem é
exato para grau 1 e perde grau 2 --- e perde \emph{muito} mais na malha adaptada
($\num{1.0e-1}$ contra $\num{3.0e-3}$), que é o resultado conhecido de que o
balanço de fluxos de segunda ordem degrada na transição 2:1. O MLS de ordem 2
recupera a exatidão de grau 2 \emph{inclusive} na malha adaptada, e o de ordem 3
é exato em grau 3 e falha em grau 4, exatamente onde deve.

Na solução trigonométrica, onde não há exatidão a afirmar, a ordem medida é a
esperada: \num{1.97} e \num{1.98} do nível 3 ao 5 para o esquema de ordem 2.
